\documentclass[11pt,letterpaper]{article}
% Shared declarations; root documents always import a concrete semantic module.
\usepackage{business-admin}
\usepackage{amsmath,amssymb}
\usepackage{longtable,tabularx,array,booktabs}
\usepackage{enumitem}
\usepackage{ragged2e}
\newcolumntype{Y}{>{\RaggedRight\arraybackslash}X}
\newcommand{\SuiteMeta}[2]{\def\DocTitle{#1}\def\DocSubtitle{#2}\def\DocVersion{1.0}\def\DocStatus{Educational reference}\def\DocOwner{Small Business Management Suite}\def\DocSummary{Original teaching and practice resource; no real-business approval implied.}}
\newcommand{\SuiteOpen}{\makedoctitle\vspace{-0.7em}\noindent\rule{\linewidth}{0.45pt}\par\medskip}
\newcommand{\SourceNote}[2]{\par\medskip\noindent{\footnotesize\textbf{Reading anchor:} Longenecker, Petty, Palich and Hoy, \emph{Small Business Management: Launching \& Growing Entrepreneurial Ventures}, 20th ed., Ch.~#1, textbook p.~#2. This document is independently authored commentary and exercises, not a reproduction of the chapter.}}
\newcommand{\WritingLines}[1]{\par\smallskip\noindent\rule{\linewidth}{0.35pt}\par\vspace{0.35cm}\noindent\rule{\linewidth}{0.35pt}\par\vspace{0.35cm}\noindent\rule{\linewidth}{0.35pt}\par\medskip}
\newcommand{\FormField}[1]{\par\medskip\noindent\textbf{#1}\par\WritingLines{3}}
\newcommand{\SBMCheckBox}{\(\square\)}
\newcommand{\AssessmentRow}[2]{#1 & #2 & \(\square\) & \(\square\) & \(\square\) & \(\square\) & \(\square\)\\\hline}
\newcommand{\AssessmentHeader}{\noindent\begin{tabular}{|p{0.22\linewidth}|p{0.49\linewidth}|c|c|c|c|c|}\hline\textbf{Criterion} & \textbf{Evidence to inspect} & \textbf{0} & \textbf{1} & \textbf{2} & \textbf{3} & \textbf{4}\\\hline}
\newcommand{\AssessmentEnd}{\end{tabular}\par}
\setlength{\emergencystretch}{3em}

\SuiteMeta{Chapter 12: Financing the Small Business}{Study guide | entrepreneurship and small business management}
\begin{document}
\SuiteOpen
\section{Learning outcome}
Compare the control, repayment, dilution and risk implications of financing alternatives.
\section{Prerequisites}
Read chapters 10, 11.\par\textbf{Required baseline:} ability to record assumptions and distinguish observations from opinions.
\section{Chapter concept map}
\begin{itemize}[leftmargin=*,itemsep=3pt]
\item debt and equity choices
\item bank underwriting
\item supplier and asset financing
\item investor capital
\item crowdfunding and public programs
\end{itemize}

\section{Technical study notes}
\begin{itemize}[leftmargin=*,itemsep=3pt]
\item Debt preserves nominal ownership but imposes mandatory payments, contractual restrictions and potentially collateral exposure.
\item Equity avoids scheduled loan repayment but changes economic participation, control and future capitalization; compare net proceeds, not headline amounts.
\item A capital request must show purpose, timing, sources, cash-service capability and contingency if financing arrives late or is smaller than expected.
\end{itemize}

\section{Reading and analysis protocol}
\begin{enumerate}[leftmargin=*,itemsep=5pt]
\item Read the chapter learning objectives, then mark terms and claims that need primary evidence.
\item For each chapter theme, distinguish the entrepreneur decision, information required, parties affected and failure condition.
\item Write a one-paragraph synthesis in your own words; cite the textbook chapter and separate contemporary practice notes.
\item Complete the applied activity and attach the underlying evidence, calculations or interview log.
\end{enumerate}

\section{Worked application and deliverable}
\textbf{Business exercise.} Prepare a lender-ready sources-and-uses and compare two realistic financing structures.\par\textbf{Required submission.} Capital strategy memorandum. Include a dated assumption register and an explicit proceed/revise/stop recommendation.\par\textbf{Assurance checks.} Debt-service headroom, ownership, collateral, documented assumptions.
\section{Common reasoning failures}
\begin{itemize}[leftmargin=*,itemsep=3pt]
\item Ignoring repayment covenants
\item Treating pledged collateral as free financing
\end{itemize}
\noindent\textbf{Correction:} require counterevidence and identify the responsible decision-maker before approving the next stage.
\section{Discussion prompt}
What evidence shows an attractive financing deal is unaffordable?\par\medskip\textbf{Evidence expected:} at least two alternatives, one quantitative or independently checkable observation, a trade-off, and a decision threshold.
\section{Exit assessment}
\begin{enumerate}[leftmargin=*,itemsep=5pt]
\item Explain each term from the concept map without copying the textbook.
\item Defend the applied decision using an assumption, evidence and counterargument.
\item Show an objective test or inspection that could prove a key assumption wrong.
\item State what would change the recommendation and who has approval authority.
\end{enumerate}

\SourceNote{12}{301}
\end{document}
